\documentclass[11pt]{article}

\usepackage[margin=25mm]{geometry}

\usepackage{amsmath,amssymb,fontspec,unicode-math}

\IfFontExistsTF{DejaVu Serif}{\setmainfont{DejaVu Serif}}{\setmainfont{Latin Modern Roman}}

\IfFontExistsTF{Latin Modern Math}{\setmathfont{Latin Modern Math}}{\setmathfont{DejaVu Math TeX Gyre}}

\usepackage{xurl,hyperref}

\setlength{\parindent}{0pt}\setlength{\parskip}{6pt}

\begin{document}

\title{Test-Time Adaptation by Causal Trimming}

\author{ELVIA preliminary AI analysis}\date{21 October 2026}\maketitle

Preliminary assessment: Minor.\par

Independent scientific review remains pending. Author adoption has not been requested.

The local trimming proposition does not automatically cover history-averaged prototypes, and nullspace removal can erase an arbitrary direction. The stated local correction result remains intact.

\section*{Review scope}

Latest-version, claim-focused review of the cited author-paper passages and applicable current author-release paths. The stated counterexamples and scope controls were checked analytically; no author program, old verifier, training run or benchmark was executed.

\section*{Problem 1: Averaged prototypes require a separate history condition}

Source: Sections 4.3–5, pp. 4–6, Eqs. (2)–(5); Appendix A.1/A.3.

Assessment: not\_issue; impact: no\_issue\_found

\subsection*{Source claim (paraphrase)}

The local trimming proposition reasons about the current boundary \(q\); deployment can instead use an averaged boundary \(\bar q\).

\subsection*{Analysis}

With \(q=(1,10,1)\), prior removed direction \((0,1,1)/\sqrt2\), current removed direction \(e_2\), and current feature \((1,-1,1)\), the local corrected margin is \(2\). The history-averaged boundary \((1,2.25,-1.75)\) instead gives margin \(-0.75\). The local conditions hold, but averaging adds an uncontrolled term.

For the deployed history rule, the signed margin is \(m_{\rm local}+y\hat z^\top(\bar q-q)\). The additional condition \(|\hat z^\top(\bar q-q)|<m_{\rm local}\) suffices for a positive margin. This explains the local theorem’s coverage; it does not amend a disproved theorem.

\subsection*{Finding and scope}

The local proposition survives; it does not alone guarantee the deployed history rule.

\section*{Problem 2: Removing beyond positive rank can erase an arbitrary null direction}

Source: Sections 4.1–4.2, p. 4, Eq. (1); Section 6, p. 7; methods/tact\_utils.py lines 214–217.

Assessment: supported; impact: minor

\subsection*{Source claim (paraphrase)}

The method removes the leading \(m\) covariance eigenvectors as augmentation-varying directions.

\subsection*{Analysis}

For covariance \(\operatorname{diag}(0,2,0)\), only \(e_2\) varies. Removing two directions permits ordered bases \((e_2,e_1,e_3)\) or \((e_2,e_3,e_1)\), which respectively erase or preserve causal feature \(e_1\). Nullspace choices do not encode observed augmentation variation.

\subsection*{Finding and scope}

The removal count needs a rank convention. This does not prove that published selected settings encounter the ambiguity or that all positive-variance directions are noncausal.

\subsection*{Proposed correction}

Cap removal by a declared numerical positive-rank threshold. If a positive-eigenvalue tie is split, specify a subspace-level or invariant tie convention.

An explicit local correction is described; author adoption and a corresponding published erratum are not established.

\section*{Limits}

Preliminary AI assessment; independent human scientific adjudication remains pending.\par

An analytic counterexample or a release discrepancy does not by itself remeasure the paper’s reported experiments.\par

\section*{Source identity}

NeurIPS 2025 official proceedings PDF; compared with current arXiv:2510.11133v1

\url{https://proceedings.neurips.cc/paper_files/paper/2025/file/09d6e2295eb50b4cdcb45d75220a2598-Paper-Conference.pdf}

Reviewed PDF SHA-256: \nolinkurl{5f486b0f762f52e5ebd5cccfe348fded3319f4bf52ea6dcbb31dc47ed09b9770}

The selected author/publisher source is the same work as the original case, matched by exact title, author identity and official version channel. The exact selected PDF SHA-256 was recomputed in the latest-version pass for Issue \#27. The PDF is a retained author/publisher copy; its original retrieval time is stated only when supported by matching acquisition metadata. File modification time was not substituted.

\end{document}
